\begin{figure}[t!]
    \captionsetup{width=1\linewidth}
	\caption{\textbf{Price response to successful vs. failed democratizations}}
	\renewcommand{\baselinestretch}{1}
    \captionsetup{singlelinecheck=off}
	\caption*{This figure presents the coefficients of the 5-year change in log prices on indicator variables for each year in a 9-year window around the end of ``successful'' and ``failed'' democratizations (Panel A) and ``liberal'' or ``reversed'' democratizations (Panel B). Successful and failed democratizations are determined using the designation in the ERT data. Namely, successful democratizations are ones in which there is a democratic transition or deepened democracy. Failed democratizations are ones in which there is no democratic transition. Liberal democratizations are ones in which the ending regime is a ``liberal democracy'' as determined by the V-Dem regime type variable. A reversed democratization is one in which the country reverts to a closed autocracy or the business or political elites become the most powerful group in the regime, also determined by the V-Dem regime indices, in the 5 years after the end of the democratization. The bars represent a 90\% confidence interval of the point estimates with standard errors clustered by year.}
	\label{fig:priceDeclineSucc}
	\footnotesize 
	\vspace{-.75cm}
    \begin{center}
		\includegraphics[width=1\textwidth]{../figures/raw/figure_D10a_democratize_price_response_succ.pdf} \\
		\includegraphics[width=1\textwidth]{../figures/raw/figure_D10b_democratize_price_response_ld.pdf}
	\end{center}
\end{figure}
